% Съдържание на български за CV
% Попълнете всеки макрос с превода на български.
% За да активирате тази версия в CI, добавете "bg" към матрицата lang в build-pdf.yml.

% Локализация
\def\country{София, България}
\def\lastupdatedtext{Последна актуализация}

% ── Образование ───────────────────────────────────────────────────────────────
\def\education{Образование}

\def\educationMSName{MSc Космическо инженерство\newline и технологии}
\def\educationMSLocation{Национален институт\newline за космически изследвания (INPE)}
\def\educationMSDuration{2018 – 2021 | Сао Жозе дос Кампос, BR}
\def\educationMSArea{Инженерство и управление на космически системи}

\def\educationBSName{BS Софтуерно инженерство}
\def\educationBSLocation{Федерален университет на Гояс (UFG)}
\def\educationBSDuration{2013 – 2016 | Гояния, BR}

% ── Умения ────────────────────────────────────────────────────────────────────
\def\skills{Умения}
\def\software{Софтуер}
\def\softwareProfessional{Професионален}
\def\softwareProjects{Проекти}
\def\softwareHobby{Хоби}

\def\skillsProfessionalList{%
Python \textbullet{} Java \textbullet{} Bash \textbullet{} SQL \\
Docker \textbullet{} Git \textbullet{} Linux \textbullet{} REST APIs%
}
\def\skillsProjectsList{%
C \textbullet{} C++ \textbullet{} ReactJS \textbullet{} R \textbullet{} \LaTeX%
}
\def\skillsHobbyList{%
Rust \textbullet{} Neovim \textbullet{} Fish Shell%
}

% ── Езици ─────────────────────────────────────────────────────────────────────
\def\languages{Езици}
\def\languagesEN{Английски: Свободно владеене — TOEFL IBT 107}
\def\languagesFR{Френски: Средно ниво}
\def\languagesJP{Японски: JLPT N5}
\def\languagesPT{Португалски: Роден}
\def\languagesBG{Български: % TODO: ниво}

% ── Професионален опит ────────────────────────────────────────────────────────
\def\professional{Професионален опит}

% EnduroSat – Технически ръководител на мисионни операции
\def\professionalESTLTitle{Технически ръководител на мисионни операции}
\def\professionalESTLDur{Февр. 2025 – Настояще | София, BG}
\def\professionalESTLDesc{%
\begin{tightemize}
\item Ръководи мисионните операции на флот от търговски спътници, определяйки технически стандарти и оперативни процедури.
\item Свързва инженерни и оперативни екипи за управление на жизнения цикъл на космическите апарати.
\item Наставничество на оператори и координация между екипи в множество паралелни мисии.
\end{tightemize}%
}

% EnduroSat – Оператор на мисия
\def\professionalESMOTitle{Оператор на мисия}
\def\professionalESMODur{Март 2023 – Февр. 2025 | София, BG}
\def\professionalESMODesc{%
\begin{tightemize}
\item Изпълняваше спътникови мисионни операции: планиране на маньоври, командване на полезни товари и диагностика на аномалии.
\item Разработваше и поддържаше оперативни наръчници и аварийни процедури за критични дейности.
\end{tightemize}%
}

% EnduroSat – Софтуерен инженер
\def\professionalESSETitle{Софтуерен инженер}
\def\professionalESSEDur{Апр. 2021 – Март 2023 | София, BG (Дистанционно)}
\def\professionalESSEDesc{%
\begin{tightemize}
\item Разработваше и поддържаше софтуер за наземния сегмент: обработка на телеметрия и планиране на мисии.
\item Изграждаше табла за мониторинг и конвейери за данни за оценка на здравето на спътника в реално време.
\end{tightemize}%
}

% По-ранен опит (накратко)
\def\professionalEarlyTitle{По-ранен опит}
\def\professionalEarlyDur{2015 – 2021}
\def\professionalEarlyDesc{%
\begin{tightemize}
\item Роли в софтуерното инженерство в MediaLab/UFG, Fitsmind (Норвегия), ZG Soluções и i9 Tecnologia — full-stack разработка на Java, Ruby on Rails, Grails и JavaScript.
\end{tightemize}%
}

% ── Научна дейност ────────────────────────────────────────────────────────────
\def\research{Научна дейност}
\def\researchMS{Докторант на пълно работно време}
\def\researchMSDur{2018 – 2021 | Сао Жозе дос Кампос}
\def\researchMSDesc{%
Разработва инструменти за Big Data анализ и алгоритми в подкрепа на операторите на спътниковия контролен център (CCS/INPE).
Стипендиант на CAPES.%
}

% ── Публикации ────────────────────────────────────────────────────────────────
\def\publications{Публикации}
